{\chapter*{Abstract}\label{sec:Abstract}}
 

\thispagestyle{plain}

\phantomsection
\addcontentsline{toc}{chapter}{Abstract}

Quantum computing offers an alternative with potential computational advantages for crucial algorithms; however, the inherent nature of quantum systems makes them highly sensitive to decoherence and environmental interactions. Consequently, developing robust error-correction strategies is a critical requirement for the advancement of the field. The continuous nature of quantum errors and the no-cloning theorem render classical redundancy approaches fundamentally unsuitable for quantum information. Nevertheless, the core principle of encoding information into a larger Hilbert space to protect it from decoherence remains valid and forms the theoretical foundation of quantum error-correcting codes. This thesis provides a mathematical introduction to quantum error correction, with a primary focus on the stabilizer formalism. First, a general quantum error model is formalized, followed by a detailed examination of the Knill-Laflamme conditions, which establish the necessary and sufficient criteria for a code to correct a specific set of errors. Furthermore, the theoretical bridge between classical linear codes and quantum stabilizer codes is developed utilizing the symplectic formalism. Finally, the practical application of this framework is demonstrated through a detailed analysis of the quantum repetition code and the 9-qubit Shor code, outlining their specific capabilities and limitations in fault-tolerant quantum computing.

\cleardoublepage
